\documentclass[a4paper,10pt]{article}
\usepackage{ctex}
\usepackage{latexsym}
\usepackage{xcolor}
\usepackage{float}
\usepackage{ragged2e}
\usepackage[empty]{fullpage}
\usepackage{wrapfig}
\usepackage{lipsum}
\usepackage{tabularx}
\usepackage{titlesec}
\usepackage{geometry}
\usepackage{marvosym}
\usepackage{verbatim}
\usepackage{enumitem}
\usepackage{hyperref}
\usepackage{fancyhdr}
\usepackage{multicol}
\usepackage{graphicx}
\usepackage{cfr-lm}
\usepackage[T1]{fontenc}
\usepackage{textcomp}
\setlength{\multicolsep}{0pt}
\pagestyle{fancy}
\fancyhf{}
\fancyfoot{}
\renewcommand{\headrulewidth}{0pt}
\renewcommand{\footrulewidth}{0pt}
\geometry{left=1.4cm, top=0.8cm, right=1.2cm, bottom=1cm}
\hypersetup{
    colorlinks=true,
    linkcolor=blue,
    urlcolor=blue,
    citecolor=blue,
}
\usepackage[most]{tcolorbox}
\tcbset{
    frame code={}
    center title,
    left=0pt,
    right=0pt,
    top=0pt,
    bottom=0pt,
    colback=gray!20,
    colframe=white,
    width=\dimexpr\textwidth\relax,
    enlarge left by=-2mm,
    boxsep=4pt,
    arc=0pt,outer arc=0pt,
}

\urlstyle{same}
\raggedright
\setlength{\tabcolsep}{0in}

\titleformat{\section}{
  \vspace{-4pt}\scshape\raggedright\large
}{}{0em}{}[\color{black}\titlerule \vspace{-7pt}]

%-------------------------
% 自定义命令
\newcommand{\resumeItem}[2]{
  \item{\textbf{#1}{:\hspace{0.5mm}#2 \vspace{-0.5mm}}}
}

\newcommand{\resumePOR}[3]{
\vspace{0.5mm}\item
    \begin{tabular*}{0.97\textwidth}[t]{l@{\extracolsep{\fill}}r}
        \textbf{#1},\hspace{0.3mm}#2 & \textit{\small{#3}}
    \end{tabular*}
    \vspace{-2mm}
}

\newcommand{\resumeProject}[4]{
\vspace{0.5mm}\item
    \begin{tabular*}{0.98\textwidth}[t]{l@{\extracolsep{\fill}}r}
        \textbf{#1} & \textit{\footnotesize{#3}} \\
        \footnotesize{\textit{#2}} & \footnotesize{#4}
    \end{tabular*}
    \vspace{-2.4mm}
}

\newcommand{\resumeSubItem}[2]{\resumeItem{#1}{#2}\vspace{-4pt}}

\renewcommand{\labelitemi}{$\vcenter{\hbox{\tiny$\bullet$}}$}

\newcommand{\resumeSubHeadingListStart}{\begin{itemize}[leftmargin=*,labelsep=0mm]}
\newcommand{\resumeHeadingSkillStart}{\begin{itemize}[leftmargin=*,itemsep=1.7mm, rightmargin=2ex]}
\newcommand{\resumeItemListStart}{\begin{justify}\begin{itemize}[leftmargin=3ex, rightmargin=2ex, noitemsep,labelsep=1.2mm,itemsep=0mm]\small}

\newcommand{\resumeSubHeadingListEnd}{\end{itemize}\vspace{2mm}}
\newcommand{\resumeHeadingSkillEnd}{\end{itemize}\vspace{-2mm}}
\newcommand{\resumeItemListEnd}{\end{itemize}\end{justify}\vspace{-2mm}}
\newcommand{\cvsection}[1]{%
\vspace{2mm}
\begin{tcolorbox}
    \textbf{\large #1}
\end{tcolorbox}
    \vspace{-4mm}
}

\newcolumntype{L}{>{\raggedright\arraybackslash}X}%
\newcolumntype{R}{>{\raggedleft\arraybackslash}X}%
\newcolumntype{C}{>{\centering\arraybackslash}X}%

%-------------------------------------------
%%%%%% 简历内容  %%%%%%%%%%%
\newcommand{\name}{王熠辉}
\newcommand{\phone}{(86) 152-6883-1452}
\newcommand{\emaila}{wangyihui@mail.hfut.edu.cn}
\newcommand{\emailb}{mistrynihilityn@gmail.com}
\newcommand{\github}{Mistry-Nihilityn}
\newcommand{\website}{https://mistry-nihilityn.github.io}

\begin{document}
\fontfamily{cmr}\selectfont

%----------头部-----------------
\parbox{3cm}{%
\includegraphics[width=2.75cm,clip]{6.png}
}\parbox{\dimexpr\linewidth-2.8cm\relax}{
\begin{tabularx}{\linewidth}{L r}
  \textbf{\LARGE \name} & \phone\\
  {软件工程 工学学士，合肥工业大学，2027届} & \href{mailto:\emaila}{\emaila} \\
  {科研实习生，新加坡国立大学 NExT++ 研究中心} & \href{mailto:\emailb}{\emailb}\\
  {浙江 杭州} & \href{https://github.com/\github}{GitHub} $|$ \href{\website}{个人主页}\\
\end{tabularx}
}

%-----------教育经历-----------------
\section{教育经历}
\setlength{\tabcolsep}{5pt}
\small{\begin{tabularx}
{\dimexpr\textwidth-3mm\relax}{|c|C|c|c|}
  \hline
  \textbf{学位} & \textbf{学校/机构} & \textbf{绩点 / 排名} & \textbf{时间}\\
  \hline
  软件工程 工学学士 & 合肥工业大学 & 90.78/100，前 1.75\% & 2023--2027\\
  \hline
  研究生（推免/保研） & 中国科学院大学 & --- & 2027--\\
  \hline
  远程科研实习 & 新加坡国立大学 NExT++ 研究中心 & --- & 2025 至今\\
  \hline
\end{tabularx}}
\vspace{-2mm}


%-----------科研经历-----------------
\section{科研经历}
\vspace{-0.4mm}
{\small\textit{全部工作于新加坡国立大学 NExT++ 研究中心完成，由杨永晖博士与蔡达成教授，NUS 讲席教授、新加坡国家科学院院士，共同指导。}}
\resumeSubHeadingListStart

    \resumeProject
      {DeepRecover：面向大语言模型的跨深度安全恢复学习}
      {第一作者，ICLR 2027 在审，CCF-A 类会议}
      {2025--2026}
      {}
      \resumeItemListStart
        \item {揭示了一个跨深度泛化差距：现有深度安全对齐方法在训练时观测到的前缀深度上能够有效恢复，但一旦评测深度超出该范围便急剧退化——问题不在于在更多位置进行恢复，而在于生成状态偏移时仍保持恢复能力。}
        \item {提出 DeepRecover：以干净查询与额外安全信息构建特权安全锚点，用共享过渡 token 纠正初始状态失配，并以逐 token 非平衡最优传输对齐轨迹失配，推理阶段无需额外干预。}
        \item {在四个指令微调大模型上仅以 30 token 前缀训练，未见深度攻击成功率降至 0.17--2.63\%，而此前深度对齐基线为 37--42\%，同时保持 MMLU 通用能力。}
    \resumeItemListEnd

    \resumeProject
      {S\textsuperscript{3}：抑制方法敏感捷径的可泛化深伪检测}
      {第一作者，目标投稿 CVPR 2027，CCF-A 类会议}
      {2025--2026}
      {}
    \resumeItemListStart
    \item {揭示深伪检测跨域失效的根因：训练后的检测器学会的是分类伪造方法而非伪造存在——这一捷径使得分布外退化在结构上不可避免，与骨干网络或数据规模无关。}
    \item {设计两种独立补救方案：梯度投影 NSP 引导训练远离方法判别性特征方向；即插即用推理模块 NAE 无需骨干重训即可达到同等效果。跨域 AUC 增益超过域内代价的 10 倍，在四个骨干网络和一个完全隔离的开放世界基准上验证。}
    \resumeItemListEnd

    \resumeProject
      {NeVA：基于神经元级编辑的 LLM 可控价值对齐}
      {共同第一作者，ACL ARR 在审，CCF-A 类会议}
      {2026 至今}
      {}
      \resumeItemListStart
        \item {定位并编辑 LLM/LVLM 中决定特定价值输出的稀疏神经元集合，无需参数更新实现细粒度价值对齐。}
    \resumeItemListEnd

  \resumeSubHeadingListEnd
\vspace{-5.5mm}


%-----------荣誉与获奖-----------------
\section{荣誉与获奖}
\vspace{-0.4mm}
\resumeSubHeadingListStart
    \resumePOR{国家奖学金，前 1\%}
        {中华人民共和国教育部}
        {2024}
    \resumePOR{金牌，CCPC 全国邀请赛}
        {南昌}
        {2026}
    \resumePOR{银牌，CCF CCSP（个人赛）}
        {成都}
        {2026}
    \resumePOR{铜牌，ICPC 亚洲区域赛}
        {成都，队伍 Hefei Uncles of Technology}
        {2026}
    \resumePOR{铜牌，ICPC 亚洲区域赛}
        {成都}
        {2025}
    \resumePOR{铜牌，2×，CCPC 区域赛}
        {济南、郑州}
        {2025}
    \resumePOR{铜牌，CCF CCSP，CSP 360/500，前 1\%}
        {}
        {2025}
    \resumePOR{一等奖，2×，安徽省机器人大赛程序设计赛道}
        {}
        {2024, 2025}
    \resumePOR{金奖，安徽省大学生创新大赛}
        {}
        {2024}
\resumeSubHeadingListEnd
\vspace{-2mm}


%-----------领导力与实践-----------------
\section{领导力与实践}
\vspace{-0.4mm}
\resumeSubHeadingListStart
    \resumePOR{ACM 队长，合肥工业大学 XCPC 校队}
        {合工大 101 队，Binary Algorithm Studio}
        {2024 至今}
    \resumeItemListStart
        \item {CCPC 全国邀请赛金牌选手；领导 20 余名队员的日常训练与题解分享。}
        \item {担任校 ACM 竞赛与安徽省程序设计竞赛命题人及组织者，服务 100 余支参赛队伍。}
    \resumeItemListEnd
    \resumePOR{视频组组长}
        {合肥工业大学全媒体中心 --- 视觉影像工作室}
        {2023--2025}
    \resumeItemListStart
        \item {制作 10 余部官方视频，多部被新华社采纳。}
        \item {组织跨学院摄影与剪辑公开课，覆盖 300 余人。}
    \resumeItemListEnd
\resumeSubHeadingListEnd
\vspace{-4mm}


%-----------技术能力-----------------
\section{技术能力}
\vspace{-0.4mm}
\resumeHeadingSkillStart
  \resumeSubItem{研究方向}
    {LLM 安全与对齐 \textbar{} 鲁棒深伪检测 \textbar{} 可解释性}
  \resumeSubItem{工程能力}
    {CCPC 全国邀请赛金牌选手；具备生产级全栈交付能力，熟练掌握 Python、C/C++、Java、Qt 及现代 Web 框架，可独立完成复杂系统从原型到部署的全流程；2 项软件著作权}
  \resumeSubItem{工具与语言}
    {PyTorch、MATLAB、\LaTeX、Git \textbar{} 英语 CET-6：553}
  \resumeSubItem{开源与社区}
    {ManimVE 渲染引擎维护者；B 站数学/CS 科普动画创作者；活跃开源贡献者}
\resumeHeadingSkillEnd
\vspace{-4mm}

%-------------------------------------------
\end{document}